\section{Логические формулы, связки и кванторы. Интерпретация, логическое следование.}

\subsection*{Логические связки и высказывания}

\paragraph{Высказывание} — это утверждение, которое либо истинно ($T$), либо ложно ($F$). Простые высказывания обозначаются пропозициональными переменными.
\paragraph{Логические связки} (в порядке старшинства):
\begin{itemize}
    \item Отрицание ($\neg$): «не».
    \item Конъюнкция ($\&$): «и».
    \item Дизъюнкция ($\lor$): «или».
    \item Импликация ($\to$): «если... то».
\end{itemize}

\subsection*{Пропозициональные формулы}

\paragraph{Синтаксис (форма Бэкуса-Наура):}
$\langle \text{формула} \rangle ::= T \mid F \mid \langle \text{переменная} \rangle \mid (\neg \langle \text{формула} \rangle) \mid (\langle \text{формула} \rangle \& \langle \text{формула} \rangle) \mid (\langle \text{формула} \rangle \lor \langle \text{формула} \rangle) \mid (\langle \text{формула} \rangle \to \langle \text{формула} \rangle)$.

\subsection*{Интерпретация и классификация формул}

\paragraph{Интерпретация $I$} — это конкретный набор истинностных значений, приписанных переменным формулы. Значение формулы $A$ в интерпретации $I$ обозначается $I(A)$.

\paragraph{Классификация формул:}
\begin{enumerate}
    \item \textbf{Выполнимая:} истинна хотя бы при одной интерпретации.
    \item \textbf{Общезначимая (тавтология):} истинна при \textbf{всех} возможных интерпретациях.
    \item \textbf{Невыполнимая (противоречие):} ложна при всех возможных интерпретациях.
\end{enumerate}

\paragraph{Теорема:} Если формулы $A$ и $A \to B$ являются тавтологиями, то формула $B$ — также тавтология (правило Modus Ponens).

\subsection*{Логическое следование и эквивалентность}

\paragraph{Логическое следование ($A \Rightarrow B$)}
Формула $B$ логически следует из $A$, если $B$ принимает значение $T$ во всех интерпретациях, где $A$ имеет значение $T$.
\begin{itemize}
    \item \textbf{Теорема 3:} $P_1 \& \dots \& P_n \Rightarrow Q$ тогда и только тогда, когда $(P_1 \& \dots \& P_n) \to Q$ — тавтология.
    \item \textbf{Теорема 4:} $P_1 \& \dots \& P_n \Rightarrow Q$ тогда и только тогда, когда $P_1 \& \dots \& P_n \& \neg Q$ — противоречие. (Это теоретическая основа метода резолюций).
\end{itemize}

\paragraph{Логическая эквивалентность ($A \equiv B$)}
Формулы логически эквивалентны, если они являются логическими следствиями друг друга (имеют одинаковые таблицы истинности). Алгебра $(\{T, F\}; \lor, \&, \neg)$ является булевой алгеброй высказываний.

\subsection*{Кванторы (Логика предикатов)}

Для работы с объектами и их свойствами в формулы вводятся \textbf{кванторы}:
\begin{itemize}
    \item \textbf{Квантор всеобщности ($\forall x$):} «для всех $x$». Формула $\forall x P(x)$ истинна, если свойство $P$ выполняется для всех объектов предметной области.
    \item \textbf{Квантор существования ($\exists x$):} «существует такой $x$, что». Формула $\exists x P(x)$ истинна, если найдется хотя бы один объект со свойством $P$.
\end{itemize}
\paragraph{Интерпретация с кванторами} включает в себя:
1. Выбор непустого множества (предметной области).
2. Определение предикатов (свойств) на этом множестве.
3. Определение значений свободных переменных.